\begin{figure}[p]\centering\begin{tikzpicture}[gclplate]\node[anchor=west,font=\sffamily\bfseries\Large,gcllabel] at (-6,3.8){UNIVERSAL CONSTRUCTION};\draw[GCLWarm,line width=1pt](-6,3.45)--(6,3.45);\node[gcllabel,draw](p) at (0,2){$f:\Pi(x:A).B(x)$};\node[gcllabel,draw](a) at (-3,.2){$a:A$};\node[gcllabel,draw](b) at (3,.2){$f\,a:B(a)$};\draw[gclbluearrow](p)--(b);\draw[gclwarmarrow](a)--(b);\end{tikzpicture}\caption{\textbf{Plate 16. Universal construction as dependent function.} A $\Pi$-term supplies evidence at each argument. Scope: the inherited $\mathrm{CHD}\text{-}1$ rules, not every semantics of universal quantification.}\label{plate:pi-universal}\end{figure}
